\section{Feature Table v1}

La representación base contiene 1,403 predictores trazables a source/family/mode. Para lectura
del equipo, las familias deben entenderse como bloques de evidencia, no como variables
intercambiables.

\subsection{Geometría y administración}

Incluye área oficial, centroides, identificadores y contexto administrativo. Estas columnas son
especialmente importantes para:
\begin{itemize}
\item joins con INEGI/SIAP;
\item cálculo de distancias geográficas;
\item control de folds agrupados;
\item análisis de soporte.
\end{itemize}

Los identificadores administrativos no se interpretan numéricamente. Cuando una variable
categórica se materializa como código, su valor entero no representa una distancia ordinal.

\subsection{BASIC histórico y 2025}

Los resúmenes BASIC se mantienen separados por sensor, índice, estadístico y ventana temporal.
Una feature típica codifica:
\[
\text{source}\times\text{period}\times\text{sensor}
\times\text{index}\times\text{statistic}\times\text{month}.
\]
Esta granularidad explica por qué la tabla base alcanza cientos de columnas aunque sólo existan
197 parcelas.

Las variables históricas 2022--2024 contribuyen al modo clean; los resúmenes 2025 completos
son competition. La distinción permite comparar:
\[
\text{estado estructural/histórico}
\quad\text{vs}\quad
\text{condición del ciclo evaluado}.
\]

\subsection{PRO Planet}

Las features Planet son competition por construcción, al corresponder a 2025. NDVI, EVI, LAI
y MSAVI se resumen de forma consistente con la entidad parcela. No se fusionan con Sentinel-2
para crear un único NDVI universal; las comparaciones inter-sensor aparecen como features
derivadas en 03A/03B.

\subsection{Clima oficial}

Se conservan precipitación, temperatura mínima y máxima históricas y 2025. Las capas
agronómicas derivan:
\[
T_{\mathrm{mean}},\quad DTR,\quad GDD_{\mathrm{proxy}},
\quad P_{\mathrm{early}},P_{\mathrm{mid}},P_{\mathrm{late}}.
\]
La existencia de la fuente base y de una transformación derivada no constituye duplicación
accidental: permite que modelos no lineales aprendan desde inputs crudos y que modelos
lineales dispongan de una base más útil.

\subsection{WaPOR}

Las 22 features de competition representan resúmenes de AETI/NPP y provenance asociado. Los
totales respetan duración de dekads. Interacciones hídricas posteriores se mantienen separadas
de estos valores source-derived.

\subsection{SoilGrids}

Las 90 features SoilGrids de la tabla clean resumen nueve propiedades a distintas
profundidades y estadísticos. Las transformaciones de 03A usan perfiles 0--30 y 30--60 para
crear gradientes y razones, evitando un polynomial expansion indiscriminado.

\subsection{SIAP}

La capa base contiene historia municipal y variables contemporáneas. Deben distinguirse:
\begin{enumerate}
\item SIAP histórico como contexto regional;
\item SIAP 2025 legado de Feature Table v1;
\item SIAP 2025 exact-scope re-auditado en 04E.1.
\end{enumerate}
La tercera es semánticamente superior para auditar la fuente, pero su RMSE directo demuestra
que mayor pureza semántica no implica mayor resolución predictiva.

\section{Agronomic Features v1}

\begin{longtable}{p{0.28\linewidth}rp{0.50\linewidth}}
\toprule
Familia & $p$ & Interpretación \\
\midrule
\endhead
phenology & 162 & forma/timing del ciclo por sensor e índice.\\
phenology anomaly & 41 & contraste 2025 contra baseline histórico parcelario.\\
water productivity & 40 & relaciones entre NPP, AETI, precipitación y crop state.\\
thermal & 36 & temperatura media, DTR y GDD proxies.\\
cross-domain & 19 & interacciones seleccionadas entre clima, suelo, SIAP y satélite.\\
soil profile & 18 & gradientes y razones verticales.\\
water timing & 10 & ventanas de precipitación y distribución estacional.\\
sensor agreement & 9 & concordancia S2--Planet en NDVI/EVI/LAI.\\
soil interaction & 8 & SOC, CEC, textura, N, pH y bases compactas.\\
nonlinear basis & 7 & log-transformaciones y otras bases target-free.\\
\bottomrule
\end{longtable}

Toda feature de esta capa puede escribirse como
\[
z_{ij}=g_j(X_i),
\]
sin dependencia de $y$. Esa propiedad permite materializarla una vez para las 197 parcelas.

\section{Empirical Features v1}

\begin{longtable}{p{0.31\linewidth}rp{0.47\linewidth}}
\toprule
Familia & $p$ & Interpretación \\
\midrule
\endhead
temporal shape & 162 & TV, roughness, curvature, autocorrelación, timing y entropía.\\
aggregate geometry & 72 & resúmenes matemáticos de bloques y distribuciones.\\
symmetric change & 42 & cambio 2025/histórico robusto a denominadores cercanos a cero.\\
short-baseline condition & 41 & posición de 2025 en rango histórico corto.\\
sensor shape similarity & 18 & Pearson/Spearman/cosine/RMSE normalizado inter-sensor.\\
\bottomrule
\end{longtable}

La capa materializada es target-free. El miner supervisado es un componente separado:
\[
\mathcal E_f=A(X_{\mathrm{train}(f)},y_{\mathrm{train}(f)}).
\]
Nunca se escribe un único set $\mathcal E$ descubierto con las 138 etiquetas completas y luego
se presenta como input independiente.

\section{Representaciones modeladas}

\subsection{Checkpoint 03C.1}

Las B-representations sirven para medir el efecto de capas:
\[
B0=\text{base},\quad
B1=\text{agro},\quad
B2=\text{empirical},
\]
\[
B3=\text{base+agro},\quad
B4=\text{base+empirical},\quad
B5=\text{all},\quad
B7=\text{all+fold-local discovery}.
\]

\subsection{Checkpoint 03C.2}

El benchmark anidado reduce el espacio:
\[
C0=\text{base},\quad
C1=\text{agro},
\]
\[
C2=\text{all deterministic + discovery},\quad
C3=\text{base+agro+selected empirical support+discovery}.
\]

\subsection{Checkpoint 04}

La topología transductiva retiene:
\[
C0,\qquad C1,\qquad C4=\text{base+agro+empirical deterministic}.
\]
C4 excluye fórmulas descubiertas con $y$ para que PCA/distancias/grafos puedan aprenderse sobre
las 197 parcelas sin contaminar la construcción $X$-only.

\section{Principio general de provenance}

Una feature derivada debe declarar:
\[
(\text{source},\text{parents},\text{formula},\text{mode},\text{evidence}).
\]
Esta tupla es más importante que el nombre humano de la columna. Dos features con nombres
parecidos pero distinto sensor, periodo o scope administrativo no deben considerarse
equivalentes.
